\section{Introduction}
\label{sec:intro}

Conditional computation is usually motivated by the same accounting argument:
most tokens do not need the whole model, so spend less on them. The argument is
made in two places, and the two answer different questions. Sparse
mixture-of-experts (MoE) layers \citep{jacobs1991moe,shazeer2017moe,fedus2022switch}
decide \emph{which} computation module processes a token, dispatching it to one
of several independent sub-networks. Recursive weight-tied architectures with
learned halting \citep{graves2016act,dehghani2019ut,banino2021pondernet} decide
\emph{how much} computation a token receives, reusing one block for a variable
number of steps.

Because the two decisions are made over different objects---which parameters
versus how many applications of them---they are orthogonal in principle, and
composing them is an obvious thing to try. Route a token to a specialised expert,
then reuse \emph{that expert's} weights for as many steps as the token appears to
need. The resulting architecture inherits an efficiency argument from each side.
What it does not inherit is evidence: whether the composition beats either
component at a matched budget has not, to our knowledge, been isolated.

We isolate it. One training engine implements three configurations differing only
in the fields their definitions require: routing without recursion (MoE, six
experts, depth 1), recursion without routing (MoR, one shared block, depth up to
7), and the composition (MoRE, six experts, depth up to 7). All three sit at
$\approx$5.58M parameters, share a bit-identical input representation, and train
on the same fixed split of WikiText-103 \citep{merity2017wikitext} with the same
optimiser budget. Because the comparison is matched rather than merely similar, a
difference between arms cannot be a difference in the harness, the data, or the
parameter count.

\paragraph{What we find.} The composition occupies the compute-efficient point of
the three, and the honest accounting of what it does and does not buy is the
paper.

\begin{itemize}[leftmargin=1.4em,itemsep=1.5pt,topsep=2pt]
\item \textbf{The composition is the most compute-efficient arm.} MoRE improves
  on MoE by $7.9\%$ validation loss ($3.948 \rightarrow 3.637$ nats/token) at
  $2.19\times$ MoE's inference FLOPs, where MoR needs $9.07\times$. It captures
  $72\%$ of MoR's quality gain over MoE for $24\%$ of MoR's inference compute---or
  $15\%$ of its \emph{extra} compute---which is roughly $3\times$ more quality per
  unit of compute than recursion alone, and $6\times$ out of domain.

\item \textbf{And the most robust out of domain.} On held-out text from a
  different distribution \citep{gao2020pile} MoRE is the \emph{best} of the three
  (7.573 nats against MoR's 7.802 and MoE's 8.272), capturing $149\%$ of MoR's
  out-of-domain gain over MoE. The in-domain ordering inverts. This is the one
  respect in which composing the two mechanisms does something neither does alone.

\item \textbf{But recursion alone remains the lowest-loss arm in domain.} MoR
  reaches 3.515 nats/token against MoRE's 3.637---the composition trades
  $0.1216$ nats, significant at our resolution floor, for a $4\times$
  inference-FLOP saving. These are two different axes and we never substitute one
  for the other: at matched parameters MoR wins on quality; per unit of inference
  compute MoRE wins. Adding routing does not improve on recursion alone in domain,
  nor does it improve routing quality over MoE alone.

\item \textbf{Recursion's adaptive depth is real; the composition's barely
  survives.} Against a \emph{depth-matched} uniform policy---forced to the learned
  policy's own mean depth, so total compute is held fixed---MoR's halting buys
  1.47 nats and MoRE's buys 0.06, a factor of 24. MoRE's halting has collapsed to
  running to its budget on 97.3\% of tokens.

\item \textbf{A diagnostic that generalises.} That depth-matched comparison is our
  most transferable contribution. A mean step count and an exit histogram---what
  this literature almost always reports---do \emph{not} establish that depth
  depends on the input; a policy ignoring its input can reproduce both. Scoring
  against a depth-matched uniform policy does, costs one inference pass, and is
  what separates our two arms, which are indistinguishable by mean-and-histogram.

\item \textbf{The two mechanisms compete for one budget.} Depth and
  specialisation are not independent knobs in MoRE the way they are in MoR and
  MoE separately: a calibration sweep that recovers adaptive depth by raising the
  ponder penalty destroys expert differentiation, and the full-corpus arm that
  preserves differentiation saturates depth instead.

\item \textbf{Two explanations are excluded, one by arithmetic.} Capacity
  dilution cannot be the cause: tokens per parameter are identical across arms by
  construction, and once recursion is counted MoRE receives slightly \emph{more}
  updates per parameter than MoR. Subspace trapping is falsified three ways---the
  recursive arms expand rather than collapse representational dimensionality,
  settle rather than oscillate, and update near-orthogonally rather than
  rescaling.
\end{itemize}

We report this as a trade-off measured at this scale rather than as a verdict on
either mechanism in general. The paper's contribution is the decomposition:
because MoE \emph{is} MoRE with the depth budget set to one and all three arms are
parameter-matched from one codebase, each contrast varies one thing. The data
scale at which the surviving account predicts the composition would overtake
recursion alone is stated as a falsifiable claim in \S\ref{sec:limitations}.

% >>>ACLONLY
\section{Related Work}
\label{sec:related}

Conditional mixtures of experts date to \citet{jacobs1991moe}; the sparsely-gated
form that made them practical is due to \citet{shazeer2017moe}, with the Top-1
dispatch and load-balancing machinery developed in GShard \citep{lepikhin2021gshard},
Switch Transformer \citep{fedus2022switch} and ST-MoE \citep{zoph2022stmoe}; their
standard justification---more parameters at constant FLOPs---is unavailable here by
construction, since we pin all three arms to one parameter count. On the depth axis,
Adaptive Computation Time \citep{graves2016act} supplies the halting unit we use,
Universal Transformers \citep{dehghani2019ut} applied it to a weight-tied block, and
PonderNet \citep{banino2021pondernet} reformulated the halting distribution to keep a
gradient near saturation; early-exit variants place the decision at the layer or token
level \citep{elbayad2020depthadaptive,schuster2022calm,raposo2024mod}, and weight tying
itself is long established \citep{lan2020albert}. That literature reports mean step
counts and exit histograms almost universally, which \S\ref{sec:halting} argues cannot
test input-dependence. Closest to us, \citet{bae2025mor} and \citet{raposo2024mod} route
\emph{to depths} over shared weights; MoRE instead routes among genuinely independent
expert FFNs whose weights the recursion then reuses, so the two decisions are made by
different heads over different objects---and neither work, to our knowledge, offers a
matched-budget comparison of such a composition against each of its components, which
is the gap we fill. We claim no novelty for the combination itself. Finally, minimal-pair
and likelihood-based probes \citep{warstadt2020blimp,ivanova2024ewok} score a
\emph{relative} preference within a matched pair, so they degrade gracefully with scale
rather than collapsing to chance, which is what makes them usable at 5.6M parameters.
% <<<ACLONLY
% >>>TMLRONLY
\section{Related Work}
\label{sec:related}

\paragraph{Sparse expert routing.}
Conditional mixtures of local experts date to \citet{jacobs1991moe}; the
sparsely-gated form that made them practical at scale is due to
\citet{shazeer2017moe}, with Top-1 dispatch, load-balancing auxiliaries and the
capacity-factor machinery developed in GShard \citep{lepikhin2021gshard}, Switch
Transformer \citep{fedus2022switch} and ST-MoE \citep{zoph2022stmoe}. The standard
justification is more parameters at roughly constant per-token FLOPs. That
argument is unavailable here by construction: we pin all three arms to the same
parameter count, so what remains of the sparsity case is a quality and robustness
case, which \S\ref{sec:results} tests directly.

\paragraph{Adaptive depth and weight-tied recursion.}
Adaptive Computation Time \citep{graves2016act} introduced the learned per-step
halting unit we use; Universal Transformers \citep{dehghani2019ut} applied it to
a weight-tied transformer block, and PonderNet \citep{banino2021pondernet}
reformulated the halting distribution to keep a gradient on the halt head near
saturation. Early-exit variants place the decision at the layer or token level
instead \citep{elbayad2020depthadaptive,schuster2022calm,raposo2024mod}. Weight
tying itself is well established as a parameter-reduction device
\citep{lan2020albert}. What this literature reports is overwhelmingly
\emph{mean} step counts and exit-rate histograms. \S\ref{sec:halting} argues that
neither statistic tests input-dependence, and that the comparison which does---against
the best uniform depth at the same budget---is rarely made.

\paragraph{Composing the two.}
Recent work has begun to combine routing with recursion, most directly by using a
router to assign per-token recursion depth \citep{bae2025mor,raposo2024mod}.
Those systems route \emph{to depths} over shared weights. MoRE differs in that
routing selects among genuinely independent expert FFNs and the recursion then
reuses \emph{that expert's} weights, so the two decisions are made by different
heads over different objects and a token keeps one expert for its whole
trajectory. We are not aware of a matched-budget comparison of such a composition
against each of its components, which is the gap this paper fills. We make no
novelty claim for the combination itself.

\paragraph{Evaluation at small scale.}
Minimal-pair benchmarks \citep{warstadt2020blimp} and likelihood-based world
knowledge probes \citep{ivanova2024ewok} score a \emph{relative} preference within a
matched pair, so they degrade gracefully with model scale rather than collapsing
to chance the way multiple-choice knowledge suites do. At 5.6M parameters they are
the appropriate instrument; \S\ref{sec:results} reports them with the caveat that
our vocabulary is fitted to the training corpus and absolute accuracies are
therefore not comparable to published models.

% <<<TMLRONLY
